\cotefile{43}
\scanpage{94}
Soit $\mathcal T'_k$ l'\unsure{espace} de Weil suivant : les objets sont les
\deleted{\unsure{quasi} \unsure{groupes} \unsure{quasi-algébriques}}
\deleted{\unsure{schémas}}
\unsure{quasi-schémas} \unsure{flasques} compacts, \unsure{la descente}
est donnée par les morphismes \deleted{\unsure{surjectifs}} \uuline{finis}
surjectifs. Soit $C'(k)$ la catégorie des faisceaux de Weil sur $\mathcal T'_k$.

\margin{Le \unsure{choix|désir} de ``fini'' est imposé par le désir d'avoir
une ``application continue'' $\mathcal T'_k\unsure{\longrightarrow|-}\mathcal T_S$ !}
\figure{\begin{tikzpicture}[x=1cm,y=1cm]
\draw(0,4.5)..controls(.9,5.1)and(.7,3.7)..(.9,3.1)..controls(1.4,1.1)and(.7,.8)..(-.1,0);
\end{tikzpicture}}

On a dans $C'(k)$ les foncteurs $H^i(k,G)$,
$\operatorname{Ext}^i(k;F,G)$, \inserted{$\underline{\operatorname{Ext}}^i_k(F,G)$}
et les relations habituelles
\[
\left\{\begin{aligned}
 H^i(k,F)&\simeq\operatorname{Ext}^i(k;\mathbb Z_k,F)\\
 \operatorname{Ext}^{\unsure{n|p+q}}(k;F,G)&\Longleftarrow
 H^p(k,\underline{\operatorname{Ext}}^{\unsure{q}}_k(F,G))
\end{aligned}\right.
\]

On notera que si $k$ est alg. clos, alors
$H^i(k,F)=0$ pour $i>0$. La formation des
$\underline{\operatorname{Ext}}^i_k(F,G)$ commute \unsure{à}
l'extension du corps de base\unsure{\ldots}

Il y a lieu de \deleted{\unsure{calculer}} donner des règles de calcul
pour des $H^i(k,G)$, $\operatorname{Ext}^i(k;F,G)$ si $G$ provient d'un
\inserted{quasi-}\unsure{schéma} en groupes habituel sur $k$,
[p.ex. d'un \unsure{pro.gr.} de Serre de type ``spécial'', i.e. à
morphismes de transition \unsure{affines|finis} pour l'indice grand\unsure{\ldots}].
Il faudrait prouver alors que ($k$ est parfait)

\noindent a) si $G$ est \unsure{simple} sur $k$, les $H^i(k,G)$ sont les
$H^i(k,\widetilde G)$, où $\widetilde G$ est le faisceau de Weil associé\ldots

\noindent b) Kif-kif pour les $\operatorname{Ext}^i(k;F,G)$\ldots
\mspage{93}

